\documentclass[a4paper,10pt,twocolumn,landscape]{article}
	\usepackage[ansinew]{inputenc}
	\usepackage{ae}
	\usepackage[brazil]{babel}
	\usepackage{xcolor}
	\usepackage{calc}
	\usepackage{wallpaper}
	\usepackage[margin=2cm]{geometry}
	\usepackage[squaren]{SIunits}
	\usepackage{amsmath}
	\usepackage{paralist}
	\usepackage{fourier}
	\usepackage{graphicx}
	\usepackage{tikz}
	\usepackage{multirow}
	\usepackage{booktabs}
	\usepackage{listings}
		
	\usetikzlibrary{calc}
		
	\graphicspath{{../figuras/}}
	
	% Minhas definições
	\newcounter{exercicio}
\newenvironment{exercicio}{%
	\refstepcounter{exercicio}%
	\penalty-200
	\noindent\colorbox{blue!60!black}{\makebox[\columnwidth-\fboxsep*2][c]{\textbf{\color{white}Exercício~\theexercicio}}}\smallskip
}{\par\medskip}
\newcommand{\foreign}[1]{\textsl{#1}}
\newcommand{\digite}[1]{{\fontfamily{cmss}\fontseries{bx}\selectfont#1}}	
\newcommand{\cs}[1]{{\normalfont\textbackslash\color{blue!50!black}#1}}
\newcommand{\pkg}[1]{{\normalfont\sffamily\color{orange}#1}}
\newcommand{\env}[1]{{\normalfont\sffamily\color{green!50!black}#1}}
\let\comando=\cs
\let\package=\pkg
\let\ambiente=\env


	% Papel de parede
	\CenterWallPaper{1}{../figuras/leao-pensador-wattermark.png}

	% Configuração extra: espaço entre as colunas.
	\setlength{\columnsep}{2cm}	
	
	
	\newsavebox{\mystarbox}
	\savebox{\mystarbox}{\color{red}$\star$}
	\newcommand{\mystar}{\usebox{\mystarbox}}
	
	\settowidth{\parindent}{\usebox{\mystarbox}}
	
		
\begin{document}

	\begin{center}\includegraphics[width=0.8\columnwidth]{../figuras/LogotipoCursoLaTeX_v2}\end{center}
	
	\section*{Figuras externas e tabelas}

	%---------------------------------------------------------
	\begin{exercicio}
		Produza cada uma das tabelas abaixo separadamente. Depois, organize-as conforme abaixo (ie, colocando-as na mesma linha e equidistantes).
		
		\medskip
		
		\noindent\hfill
		\begin{tabular}{cc}		
			1 & 2 \\		
			3 & 4 \\
		\end{tabular}\hfill
		\begin{tabular}{|c|c|}
			1 & 2 \\ 
			3 & 4 \\ 
		\end{tabular}\hfill
		\begin{tabular}{cc}
			\hline
			1 & 2 \\ 
			\hline
			3 & 4 \\
			\hline
		\end{tabular}\hfill
		\begin{tabular}{|c|c|}
			\hline
			1 & 2 \\ 
			\hline
			3 & 4 \\ 
			\hline
		\end{tabular}\hfill\null
		
		\medskip
		
		\footnotesize
		\noindent\textbf{obs.:} use \cs{hfill} entre as tabelas e, após a última, 
		use \cs{hfill}\cs{null} ou \cs{hfill\{\}}. Isto é necessário porque o 
		\LaTeX\ ignora espaços horizontais no final da linha (ou espaços verticais 
		no final da página).
	\end{exercicio}
	
	%---------------------------------------------------------
	\begin{exercicio}
		Reproduza a equação matricial abaixo (\textbf{atenção:} ``in'' e ``out'' são palavras, não variáveis).
		
		\medskip
	
		\[
		\left(\begin{array}{cc}
			E - H_C - \Sigma_T & i\Sigma^\text{in} \\
			i\Sigma^\text{out} & E - H_C + \Sigma^{\bar T}
		\end{array}\right)
		\left(\begin{array}{cc}
			G^T & -iG^n \\
			-iG^p & -G^{\bar T}
		\end{array}\right) = 
		\left(\begin{array}{cc}
			I & 0 \\ 0 & I
		\end{array}\right)
		\]
	\end{exercicio}
	
	%---------------------------------------------------------
	\begin{exercicio}
		Utilize o comando \cs{cline} para produzir estas tabelas:
	
		\medskip
		
		\noindent\hfill
		\begin{tabular}{ccc}
			\cline{1-1}
			1 & 2 & 3 \\
			\cline{2-2}
			4 & 5 & 6 \\
			\cline{3-3}
		\end{tabular}\hfill
		\begin{tabular}{ccc}
			\cline{1-1}\cline{3-3}
			1 & 2 & 3 \\
			\cline{2-2}
			4 & 5 & 6 \\
			\cline{1-3}
		\end{tabular}\hfill{}
	\end{exercicio}
		
	%---------------------------------------------------------
	\begin{exercicio}
		Utilize o comando \cs{multicolumn} para produzir as tabelas abaixo.
		
		\medskip		
		
		\noindent\hfill
		\begin{tabular}{|c|c|c|}
			\hline
			1 & 2 & 3 \\
			\hline
			4 & 5 & 6 \\
			\hline
		\end{tabular}\hfill
		\begin{tabular}{|c|c|c|}
			\hline
			\multicolumn{2}{|c|}{1 e 2} & 3 \\
			\hline
			4 & \multicolumn{2}{c|}{5 e 6} \\
			\hline
		\end{tabular}\hfill
		\begin{tabular}{|c|c|c|}
			\hline
			\multicolumn{3}{|c|}{1, 2 e 3} \\
			\hline
			\multicolumn{3}{|c|}{4, 5 e 6} \\
			\hline
		\end{tabular}\hfill{}
		
	\end{exercicio}

	%---------------------------------------------------------
	\begin{exercicio}
		Utilize o comando \cs{multirow}, definido no pacote \pkg{multirow}, para 
		produzir estas tabelas (as células unificadas têm \unit{1}{\centi\metre} de
		largura):
		
		\medskip
		
		\noindent\hfill
		\begin{tabular}{|c|c|c|}
			\hline
			1 & 2 & 3 \\
			\hline
			4 & 5 & 6 \\
			\hline			
		\end{tabular}\hfill
		\begin{tabular}{|c|c|c|}
			\hline
			\multirow{2}{1cm}{\centering 1 e 4} & 2 & 3 \\
			                         & 5 & 6 \\
			\hline			
		\end{tabular}\hfill
		\begin{tabular}{|c|c|c|}
			\hline
			1 & \multirow{2}{1cm}{\centering 2 e 5} & 3 \\
			4 &                          & 6 \\
			\hline			
		\end{tabular}\hfill
		\begin{tabular}{|c|c|c|}
			\hline
			1 & 2 & \multirow{2}{1cm}{\centering 3 e 6} \\
			4 & 5 &                          \\
			\hline			
		\end{tabular}\hfill{}
		
		\medskip
		
		\footnotesize
		\noindent\textbf{dica:} utilize \cs{centering} para centralizar o
		conteúdo das células unificadas.
	\end{exercicio}
		
	%---------------------------------------------------------
	\begin{exercicio}
		Utilize o pacote \pkg{booktabs} para produzir a seguinte tabela. Isto é, 
		utilize o comando \cs{toprule} para a linha horizontal superior,
		\cs{cmidrule} e \cs{midrule} para as linhas no meio da tabela, e
		\cs{bottomrule} para terminá-la.
		
		\medskip
		
		\centering
		\begin{tabular}{llr}			
			\toprule
			\multicolumn{2}{c}{Item} \\
			\cmidrule{1-2}
			Animal & Description & Price (\$) \\
			\midrule
			Gnat      & per gram & 13.65 \\
			          & each     &  0.01 \\
			Gnu       & stuffed  & 92.50 \\
			Emu       & stuffed  & 33.33 \\
			Armadillo & frozen   & 8.99  \\
			\bottomrule
		\end{tabular}
	\end{exercicio}

	%---------------------------------------------------------
	\begin{exercicio}
		Refaça o exercício anterior sem utilizar o pacote \pkg{booktabs}, ou seja,
		utilizando apenas os comandos \cs{hline} e \cs{cline}. Observe as 
		diferenças.
	\end{exercicio}

	%---------------------------------------------------------
	\begin{exercicio}
		Crie um documento \LaTeX\ e, utilizando o pacote \pkg{graphicx}, importe 
		para dentro dele uma figura no formato \texttt{eps} com altura igual a 
		$1/5$ da altura do texto (comprimento \cs{textheight}). Para qual(is) 
		formato(s) de documento de saída você pode inserir uma figura deste tipo?
	\end{exercicio}
		
	%---------------------------------------------------------
	\begin{exercicio}
		Se você quiser incluir figuras \texttt{pdf}, \texttt{jpg} e/ou 
		\texttt{png} no seu	documento	\LaTeX, qual deve ser o formato do documento
		de saída?
	\end{exercicio}
	
	\newsavebox{\minhacaixa}
	\savebox{\minhacaixa}{%
		\includegraphics[width=0.09\columnwidth]{Tei-Gi}}
	
	%---------------------------------------------------------
	\begin{exercicio}
		Coloque dez figuras uma do lado da outra, numa mesma linha, \emph{sem 
		ultrapassar as margens} (isto é, sem produzir \foreign{bad boxes}):
		
		\medskip
		
		\noindent
		\hfill\usebox{\minhacaixa}%
		\hfill\usebox{\minhacaixa}%
		\hfill\usebox{\minhacaixa}%
		\hfill\usebox{\minhacaixa}%
		\hfill\usebox{\minhacaixa}%
		\hfill\usebox{\minhacaixa}%
		\hfill\usebox{\minhacaixa}%
		\hfill\usebox{\minhacaixa}%
		\hfill\usebox{\minhacaixa}%
		\hfill\usebox{\minhacaixa}%
		\hfill\null
	\end{exercicio}
		
	%---------------------------------------------------------	
	\savebox{\minhacaixa}{%
		\setlength{\fboxsep}{0pt}
		\fbox{\includegraphics[width=0.09\columnwidth]{Tei-Gi}}}
						
	\begin{exercicio}
		Refaça o exercício anterior, mas coloque também uma moldura ao redor da 
		figura:
		
		\medskip
		
		\noindent
		\hfill\usebox{\minhacaixa}%
		\hfill\usebox{\minhacaixa}%
		\hfill\usebox{\minhacaixa}%
		\hfill\usebox{\minhacaixa}%
		\hfill\usebox{\minhacaixa}%
		\hfill\usebox{\minhacaixa}%
		\hfill\usebox{\minhacaixa}%
		\hfill\usebox{\minhacaixa}%
		\hfill\usebox{\minhacaixa}%
		\hfill\usebox{\minhacaixa}%
		\hfill\null
		
		\medskip
		
		\footnotesize
		\noindent\textbf{obs.:} torne nulo o comprimento \cs{fboxsep} para 
		facilitar.
	\end{exercicio}

	%---------------------------------------------------------
	\begin{exercicio}
		Utilize apenas o comando \cs{includegraphics} para inserir uma figura 
		rotacionada de \unit{90}{\degree}:
		
		\medskip
		
		\centering
		\includegraphics[width=2cm,angle=90]{rosa_dos_ventos}
	\end{exercicio}
	
	%---------------------------------------------------------
	\begin{exercicio}
		Utilize os comandos \cs{raisebox} e \cs{includegraphics} para produzir o 
		resultado a seguir (a linha horizontal representa a linha-base):
		
		\medskip
		
		\noindent
		\hrulefill\raisebox{+0.0\height}{\includegraphics[width=1cm]{Tei-Gi}}%
		\hrulefill\raisebox{-0.5\height}{\includegraphics[width=1cm]{Tei-Gi}}%
		\hrulefill\raisebox{-1.0\height}{\includegraphics[width=1cm]{Tei-Gi}}%
		\hrulefill\null
		
		\medskip
		
		\footnotesize
		\noindent\textbf{dica:} quando você cria uma caixa LR com os comandos
		\cs{raisebox}, \cs{framebox} ou \cs{makebox}, o \LaTeX\ define
		automaticamente os seguintes comprimentos: \cs{width}, \cs{height}, 
		\cs{depth} e \cs{totalheight}, que representam a largura, altura,
		profundidade e altura-total \emph{natural} desta caixa,	respectivamente.
		Mas atenção: esses comprimentos só valem dentro dos	argumentos (opcionais
		ou obrigatórios) do comando. Por exemplo, eu posso definir	uma caixa LR
		de comprimento igual ao dobro do seu comprimento natural assim: 
		\cs{framebox}\verb|[2|\cs{width}\verb|]{|\textit{\textbf{caixa}}\verb|}|.
	\end{exercicio}

	%---------------------------------------------------------
	\begin{exercicio}		
		Utilize o comando \cs{resizebox} para produzir a linha principal do 
		exercício $6$ dos exercícios sobre caixas, que ocupa precisamente toda a linha do texto:
		
		\medskip
		
		\noindent\resizebox{\columnwidth}{!}{Identifique a linha-base.}
		
		\medskip
	
		\footnotesize
		\noindent\textbf{obs.:} não se esqueça de suprimir a endentação.
	\end{exercicio}
	
	%---------------------------------------------------------
	\begin{exercicio}
		Reveja os exercícios 11 e 12 da lista anterior e produza o resultado 
		abaixo. Tente	responder aqueles exercícios agora.
		
		\Huge
		\noindent\hrulefill\scalebox{1}[-1]{Eu sou uma caixa!}\hrulefill
	\end{exercicio}
	
	%---------------------------------------------------------
	\begin{exercicio}
		Utilize o comando \cs{rotatebox} para produzir
		
		\medskip
	
		\centering
		\rotatebox{155}{%
		\begin{tabular}{|c|c|c|}
			\hline
			1 & 2 & 3 \\\hline 4 & 5 & 6\\
			\hline
		\end{tabular}}
	\end{exercicio}
	
	%---------------------------------------------------------
	\begin{exercicio}
		Utilize a mesma tabela do exercício anterior junto com o comando
		\cs{resizebox} e reproduza o resultado abaixo, cuja largura é metade da
		da linha e a altura, \unit{1}{\centi\metre}.
		
		\medskip
	
		\centering
		\resizebox{0.5\columnwidth}{1cm}{%
		\begin{tabular}{|c|c|c|}
			\hline
			1 & 2 & 3 \\\hline 4 & 5 & 6\\
			\hline
		\end{tabular}}
	\end{exercicio}
		
	%---------------------------------------------------------
	\begin{exercicio}
		\textbf{Desafio:} construa a célula unificada da tabela abaixo sem
		utilizar o pacote \pkg{multirow}.

		\setlength{\fboxsep}{0.5pt}

		\medskip

		\begin{center}
			\begin{tabular}{|c|c|}		
				\hline
				Célula 11 &\\
				\cline{1-1}
				Célula 21 & \rule{3mm}{0pt}%
					\raisebox{1ex}[\height][0pt]{União}%
					\rule{3mm}{0pt} \\
				\hline
			\end{tabular}
		\end{center}
		
		\medskip
		
		\footnotesize
		\noindent\textbf{dica:} a célula unificada em questão não é, de fato,
	unificada. Acontece que a célula $(\text{linha},\text{coluna}) = (1,2)$ está 
	vazia, enquanto $(2,2)$ contém um \cs{raisebox}. Experimente.

	\end{exercicio}

		
\end{document}